% ============================================================
% chapter5.tex  —  Chapitre 5 : Discussion, Limitations
%                  and Future Perspectives
% ============================================================

\chapter{Discussion, Limitations and Future Perspectives}


\noindent
The previous chapters described the design, architecture,
and implementation of DecisioBI.
This final chapter reflects on the results obtained,
discusses the limitations encountered during development,
and outlines perspectives for future improvements
and extensions.


% ============================================================
\section{Discussion}
% ============================================================

\noindent
DecisioBI demonstrates that a modular, open-source
Business Intelligence platform can be built
to address the specific needs of SMEs
in resource-constrained environments.
The following discussion examines
how the platform meets its original objectives
and where trade-offs were necessary.

% --- 5.1.1 ---
\subsection{Modularity and Separation of Concerns}

\noindent
The decision to decompose the platform into nine
independent Django applications proved effective.
Each module (\texttt{ingestion}, \texttt{nettoyage},
\texttt{conflits}, \texttt{kpi}, \texttt{dashboard},
\texttt{ia\_interpretation}, \texttt{chatbot},
\texttt{notifications}, \texttt{authentication})
maintains its own models, services, serializers,
and API endpoints.
This separation made it possible to develop
and test each module independently,
reducing integration risk
and simplifying debugging.

\vspace{0.3cm}
\noindent
However, this modularity introduced a trade-off:
some degree of code duplication exists
across modules, particularly in serialisation
logic and error handling.
A shared utility layer could reduce this duplication
without compromising the separation of responsibilities.

% --- 5.1.2 ---
\subsection{Semi-Automatic Pipeline Design}

\noindent
One of the most deliberate architectural choices
was making the processing pipeline semi-automatic
rather than fully automated.
After uploading a file, the system performs
ingestion and schema profiling automatically,
but the user must explicitly trigger
and validate each subsequent stage —
cleaning, conflict detection, and KPI calculation.

\vspace{0.3cm}
\noindent
This human-in-the-loop approach was motivated
by the nature of SME data.
Operational data from small businesses
is often messy, inconsistent,
and context-dependent.
Fully automated pipelines risk propagating
errors silently.
By requiring user validation at critical decision points,
DecisioBI ensures that business users
retain control over their data
while benefiting from automated analysis.

\vspace{0.3cm}
\noindent
The downside is that this approach
increases the time required to go
from raw data to actionable dashboard.
For organisations with large volumes of data,
a more automated pipeline with optional override
might be preferable.

% --- 5.1.3 ---
\subsection{AI Integration}

\noindent
The integration of the Groq API
with the Llama~3.3~70B model
provides natural-language interpretation
of KPI data and a conversational chatbot
interface.
This approach offers several advantages:
no model hosting is required,
the API is fast and reliable,
and the model produces coherent analyses
in both French and English.

\vspace{0.3cm}
\noindent
The main limitation is dependency on an external service.
If the Groq API becomes unavailable
or changes its pricing model,
both the KPI interpretation and chatbot features
become non-functional.
A fallback mechanism or local model
could mitigate this risk.

% --- 5.1.4 ---
\subsection{Testing and Quality Assurance}

\noindent
The test suite comprises over 478 test functions
and classes across 35 test files,
covering model constraints, API behaviour,
service logic, and security enforcement.
The 70\% coverage target was set
as a pragmatic minimum —
high enough to catch critical defects
without imposing an unrealistic burden
on a small development team.

\vspace{0.3cm}
\noindent
The testing strategy validated
that RBAC enforcement, data isolation,
and input validation work as specified.
However, performance testing
and load testing were not part
of the evaluation scope,
which represents a gap
that should be addressed
before production deployment.


% ============================================================
\section{Limitations}
% ============================================================

\noindent
Despite the progress achieved,
DecisioBI has several limitations
that should be acknowledged honestly.

% --- 5.2.1 ---
\subsection{Deployment and Infrastructure}

\noindent
The platform was developed and tested
in a local Windows environment.
Several infrastructure components
are not fully operational:

\begin{itemize}
  \item \textbf{Redis is not installed} on the development
        machine, which means Celery workers
        and background task processing
        cannot be tested locally.
        All asynchronous operations
        (file ingestion, data cleaning,
        KPI calculation) require Redis
        to function.
  \item \textbf{No containerisation} (Docker)
        or continuous integration / continuous
        deployment (CI/CD) pipeline exists.
        Deployment to production
        would require manual configuration.
  \item \textbf{No reverse proxy} or SSL termination
        is configured for the backend.
        Production deployment would require
        nginx or a similar solution.
\end{itemize}

% --- 5.2.2 ---
\subsection{Data Processing}

\noindent
The cleaning module provides 18 rule types
and a pipeline engine,
but several limitations remain:

\begin{itemize}
  \item \textbf{No data versioning:}
        once data is cleaned,
        the original raw data
        cannot be recovered
        through the user interface.
        The \texttt{RawData} table
        preserves the original upload,
        but no rollback mechanism exists.
  \item \textbf{Semi-manual workflow:}
        as discussed in Section~5.1.2,
        the requirement for user validation
        at each stage
        limits the platform's ability
        to process data autonomously.
  \item \textbf{No streaming ingestion:}
        the platform only supports
        file-based upload.
        Real-time data streams
        (e.g., from APIs or message queues)
        are not supported.
\end{itemize}

% --- 5.2.3 ---
\subsection{Analytics and Intelligence}

\noindent
The KPI module provides variance analysis,
Z-score anomaly detection,
and linear regression forecasting.
However:

\begin{itemize}
  \item \textbf{No multivariate anomaly detection:}
        the Z-score method evaluates
        each KPI independently.
        The Isolation Forest algorithm
        was implemented in a separate module
        but is not currently exposed
        in the user interface.
  \item \textbf{No predictive analytics beyond linear models:}
        the forecasting capability
        is limited to linear regression.
        Seasonal patterns, exponential smoothing,
        and machine learning-based forecasting
        are not implemented.
  \item \textbf{No alerting system:}
        while the \texttt{KPIAlert} model exists,
        the multi-channel notification system
        (email, webhook, Slack)
        is not fully integrated
        into the user workflow.
\end{itemize}

% --- 5.2.4 ---
\subsection{User Experience}

\noindent
The frontend provides a functional interface,
but several UX limitations exist:

\begin{itemize}
  \item \textbf{No mobile optimisation:}
        the interface is designed
        for desktop browsers.
        Responsive design is partially implemented
        through Tailwind CSS,
        but no dedicated mobile experience exists.
  \item \textbf{French-only interface:}
        while the chatbot supports bilingual
        interaction (French/English),
        the user interface itself
        is only available in French.
  \item \textbf{No offline capability:}
        the platform requires
        a persistent internet connection.
        No service worker or caching strategy
        is implemented for offline access.
  \item \textbf{No data export beyond PDF:}
        users can export dashboards to PDF,
        but CSV or Excel export
        of dashboard data is not available
        from the dashboard page.
\end{itemize}

% --- 5.2.5 ---
\subsection{Security}

\noindent
The security implementation covers
authentication, authorisation, and data isolation.
However:

\begin{itemize}
  \item \textbf{No rate limiting:}
        beyond Django REST Framework's
        default throttling,
        no rate limiting is configured
        for API endpoints.
        This could leave the platform
        vulnerable to brute-force attacks.
  \item \textbf{No audit logging:}
        while the \texttt{conflits} module
        includes an \texttt{ActivityLog} model,
        no system-wide audit trail
        exists for tracking user actions
        across all modules.
  \item \textbf{No HTTPS enforcement:}
        the development server
        runs over HTTP.
        Production deployment
        would require TLS configuration.
\end{itemize}

% --- 5.2.6 ---
\subsection{Testing Scope}

\noindent
While the test suite is substantial,
several areas are not covered:

\begin{itemize}
  \item \textbf{No performance or load testing:}
        response times under heavy load
        are not measured.
  \item \textbf{No end-to-end testing:}
        no browser automation tests
        (e.g., Cypress, Playwright)
        validate the complete user workflow.
  \item \textbf{No cross-browser testing:}
        the frontend was tested
        primarily on Chromium-based browsers.
\end{itemize}


% ============================================================
\section{Future Perspectives}
% ============================================================

\noindent
The limitations identified above
suggest several directions
for future development.

% --- 5.3.1 ---
\subsection{Infrastructure Improvements}

\noindent
The highest priority improvements
are infrastructure-related:

\begin{enumerate}
  \item \textbf{Containerisation:}
        Docker Compose should be configured
        to package the backend, frontend,
        PostgreSQL, and Redis
        into a single deployable unit.
        This would eliminate the dependency
        on manual server configuration
        and enable reproducible deployments.
  \item \textbf{CI/CD pipeline:}
        a GitHub Actions workflow
        should be implemented
        to run tests automatically
        on each commit,
        build the frontend,
        and deploy to production.
  \item \textbf{Production-ready deployment:}
        nginx as a reverse proxy,
        Gunicorn as the WSGI server,
        and Let's Encrypt for TLS
        would make the platform
        production-ready.
\end{enumerate}

% --- 5.3.2 ---
\subsection{Data Processing Enhancements}

\noindent
Several improvements would strengthen
the data processing pipeline:

\begin{enumerate}
  \item \textbf{Data versioning:}
        implementing a versioning system
        for cleaned data
        would allow users to compare
        different cleaning iterations
        and roll back if necessary.
  \item \textbf{Automated cleaning mode:}
        an optional fully automated mode
        would allow advanced users
        to skip manual validation
        when processing trusted data sources.
  \item \textbf{Streaming ingestion:}
        support for API-based
        and message queue-based data sources
        would extend the platform's applicability
        beyond file-based imports.
  \item \textbf{Multi-sheet Excel support:}
        allowing users to select
        specific sheets from Excel files
        and define inter-sheet relationships
        would improve compatibility
        with real-world ERP exports.
\end{enumerate}

% --- 5.3.3 ---
\subsection{Analytics and Intelligence}

\noindent
The analytical capabilities
could be extended in several ways:

\begin{enumerate}
  \item \textbf{Advanced forecasting:}
        exponential smoothing,
        ARIMA models,
        and Prophet-based forecasting
        would provide more accurate predictions
        for seasonal and non-linear data.
  \item \textbf{Multivariate anomaly detection:}
        re-exposing the Isolation Forest algorithm
        through the user interface
        would enable detection of complex
        multi-dimensional outliers.
  \item \textbf{Custom KPI formulas:}
        allowing users to define KPIs
        through a visual formula builder
        rather than writing SQL or Python
        would lower the barrier to entry.
  \item \textbf{Drill-down capabilities:}
        enabling users to navigate
        from aggregate KPI values
        to underlying raw data
        would improve analytical depth.
\end{enumerate}

% --- 5.3.4 ---
\subsection{User Experience Improvements}

\noindent
The user experience could be enhanced through:

\begin{enumerate}
  \item \textbf{Responsive mobile design:}
        optimising the interface
        for tablet and smartphone access
        would make the platform accessible
        to users who primarily work on mobile devices.
  \item \textbf{Internationalisation:}
        adding English and local language
        support for the user interface
        would broaden the platform's reach.
  \item \textbf{Onboarding workflow:}
        a guided tutorial
        for new users
        would reduce the learning curve
        and improve adoption rates.
  \item \textbf{Real-time notifications:}
        WebSocket-based push notifications
        would replace the current polling mechanism
        and provide a more responsive experience.
\end{enumerate}

% --- 5.3.5 ---
\subsection{SaaS and Multi-Tenancy}

\noindent
The current architecture supports
organisation-based data isolation
through foreign key scoping,
but lacks several features
required for a true SaaS deployment:

\begin{enumerate}
  \item \textbf{Subscription and billing management:}
        a Stripe or LemonSqueezy integration
        would enable tiered pricing plans.
  \item \textbf{Usage quotas:}
        enforcing limits on data sources,
        KPIs, and API calls per organisation
        would prevent resource abuse.
  \item \textbf{Tenant administration:}
        self-service organisation creation,
        invitation workflows,
        and role management
        would reduce operational overhead.
\end{enumerate}

% --- 5.3.6 ---
\subsection{Reporting and Compliance}

\noindent
The notification module provides
PDF report generation,
but a complete reporting system would include:

\begin{enumerate}
  \item \textbf{Scheduled reports:}
        automatic generation and email delivery
        of daily, weekly, or monthly reports.
  \item \textbf{Excel export:}
        full-fidelity export of dashboards
        and pivot tables to Excel format.
  \item \textbf{Audit trail:}
        system-wide logging of user actions
        for regulatory compliance.
  \item \textbf{Data lineage:}
        tracking the origin and transformation
        history of each data point
        from ingestion to dashboard.
\end{enumerate}


% ============================================================
\section{Chapter Summary}
% ============================================================

\noindent
DecisioBI demonstrates that a modular,
open-source Business Intelligence platform
can be built to address SME needs
in resource-constrained environments.
The iterative development approach,
the semi-automatic pipeline design,
and the integration of large language models
for data interpretation
represent practical engineering decisions
shaped by the realities of small-team development.

\vspace{0.3cm}
\noindent
The platform's main strengths
are its modular architecture,
its human-centred pipeline design,
and its accessible AI integration.
Its main limitations
are infrastructure readiness,
the lack of advanced analytics,
and the absence of mobile optimisation.

\vspace{0.3cm}
\noindent
Future work should prioritise
containerisation and production deployment,
advanced forecasting and anomaly detection,
and the extension of the platform
into a multi-tenant SaaS offering.
These improvements would transform DecisioBI
from a thesis prototype
into a deployable product
capable of serving SMEs across West Africa.

\vspace{0.7cm}
\begin{center}
  \textcolor{BITred!40}{\rule{0.35\linewidth}{0.5pt}}
  \hspace{0.5em}
  \textcolor{BITred}{$\star$}
  \hspace{0.5em}
  \textcolor{BITred!40}{\rule{0.35\linewidth}{0.5pt}}
\end{center}
